\documentclass[11pt]{article}
\usepackage{amsmath,amssymb,amsthm}
\usepackage{algorithm,algorithmic}
\usepackage{graphicx}
\usepackage{hyperref}
\usepackage{cleveref}
\usepackage[margin=1in]{geometry}
\usepackage{xcolor}

\newtheorem{theorem}{Theorem}
\newtheorem{lemma}{Lemma}
\newtheorem{proposition}{Proposition}
\newtheorem{definition}{Definition}
\newtheorem{remark}{Remark}

\title{Deep Learning for Bivariate Copula Density Estimation:\\
Why Full Diffusion Models Fail and What Actually Works for Vine Copulas}
\author{}
\date{November 2025}

\begin{document}

\maketitle

\begin{abstract}
We investigate the application of deep learning to bivariate copula density estimation for vine copula construction. Our primary goal is \emph{amortized pair-copula fitting}: given empirical data for each pair, rapidly produce valid copula densities and $h$-functions for vine recursion. We document a failed attempt using unconditional denoising diffusion probabilistic models (DDPMs) in log-density space, which catastrophically collapsed to uniform predictions (MAE $\approx 1.0$). Through careful analysis, we identify the root causes: (1) objective mismatch between log-space diffusion and density-space constraints, (2) incorrect likelihood formulation favoring uniform densities, and (3) post-hoc projection that flattens unstructured predictions. We propose a practical solution: a \textbf{noise-conditioned density denoiser} that treats diffusion as a training augmentation strategy rather than a generative sampling process. This approach provides single-pass inference with exact copula constraints via Sinkhorn projection, suitable for efficient vine construction. We also discuss full generative modeling approaches (latent diffusion, score-based models) as separate research directions for unconditional copula generation.
\end{abstract}

\section{Introduction}

\subsection{Problem Statement}

A bivariate copula is a function $C: [0,1]^2 \to [0,1]$ that satisfies:
\begin{enumerate}
    \item \textbf{Grounded margins}: $C(u, 0) = C(0, v) = 0$ for all $u,v \in [0,1]$
    \item \textbf{Uniform margins}: $C(u, 1) = u$ and $C(1, v) = v$ for all $u,v \in [0,1]$
    \item \textbf{2-increasing}: For all $u_1 \leq u_2$ and $v_1 \leq v_2$:
    \begin{equation}
    C(u_2, v_2) - C(u_2, v_1) - C(u_1, v_2) + C(u_1, v_1) \geq 0
    \end{equation}
\end{enumerate}

The copula density $c(u,v) = \frac{\partial^2 C}{\partial u \partial v}$ must satisfy:
\begin{align}
c(u,v) &\geq 0 \quad \forall (u,v) \in [0,1]^2 \label{eq:positivity} \\
\int_0^1 \int_0^1 c(u,v) \, du \, dv &= 1 \label{eq:normalization} \\
\int_0^1 c(u,v) \, dv &= 1 \quad \forall u \in [0,1] \label{eq:uniform_u} \\
\int_0^1 c(u,v) \, du &= 1 \quad \forall v \in [0,1] \label{eq:uniform_v}
\end{align}

The goal is to learn a generative model $p_\theta(c)$ that produces valid copula densities.

\subsection{Motivation for Generative Modeling}

Traditional parametric copula families (Gaussian, Clayton, Gumbel, etc.) are limited in expressiveness. Non-parametric methods (kernel density estimation, Bernstein copulas) require large sample sizes. Deep generative models offer a promising middle ground, potentially learning flexible copula families from data while maintaining structural constraints.

\subsection{Two Distinct Use Cases}

It is critical to distinguish two different objectives:

\paragraph{Use Case 1: Unconditional Copula Generation}
Learn a generative model $p_\theta(c)$ that can sample diverse, novel copula densities from a learned distribution. This is analogous to generative image modeling (GANs, diffusion models for images). Applications include:
\begin{itemize}
    \item Exploring the space of possible dependence structures
    \item Simulation studies requiring synthetic copula families
    \item Learning copula priors from large datasets
\end{itemize}

\paragraph{Use Case 2: Amortized Pair-Copula Estimation for Vines}
Given empirical pseudo-observations $(u_i, v_i)$ for a specific pair of variables, rapidly produce:
\begin{itemize}
    \item A valid copula density $\hat{c}(u,v)$ fitting the data
    \item The corresponding $h$-functions $h(u|v)$ and $h(v|u)$ for vine recursion
\end{itemize}
This is a \textbf{conditional density regression} problem, not unconditional generation. Requirements:
\begin{itemize}
    \item Fast single-pass inference (thousands of pairs in a vine)
    \item Robust to varying sample sizes and noise
    \item Exact satisfaction of copula constraints
\end{itemize}

\textbf{This paper primarily addresses Use Case 2.} We initially attempted a full diffusion approach (Use Case 1) and document why it failed for our needs. The recommended solution treats diffusion as a \emph{training augmentation technique}, not as a multi-step generative sampler.

\section{Negative Result: Unconditional Log-Space DDPM}

\subsection{Architecture Overview}

The current implementation employs a denoising diffusion probabilistic model (DDPM) \cite{ho2020denoising} that operates on the \textbf{log-density} representation $\ell(u,v) = \log c(u,v)$.

\subsubsection{Forward Diffusion Process}

Given a copula density $c(u,v)$, we compute its log-density:
\begin{equation}
\ell_0(u,v) = \log c(u,v)
\end{equation}
clipped to the range $[-20, 20]$ to prevent numerical overflow.

The forward process adds Gaussian noise over $T=400$ timesteps:
\begin{equation}
q(\ell_t | \ell_0) = \mathcal{N}(\ell_t; \sqrt{\bar{\alpha}_t} \ell_0, (1-\bar{\alpha}_t) \mathbf{I})
\end{equation}
where $\bar{\alpha}_t = \prod_{s=1}^t \alpha_s$ with a cosine noise schedule:
\begin{equation}
\bar{\alpha}_t = \cos^2\left(\frac{t/T + s}{1+s} \cdot \frac{\pi}{2}\right), \quad s = 0.008
\end{equation}

\subsubsection{Reverse Denoising Process}

A U-Net model $\epsilon_\theta(\ell_t, t)$ predicts the noise added at timestep $t$:
\begin{equation}
\mathcal{L}_{\text{noise}} = \mathbb{E}_{t, \ell_0, \epsilon} \left[ \|\epsilon - \epsilon_\theta(\ell_t, t)\|^2 \right]
\end{equation}

To reconstruct the clean log-density:
\begin{equation}
\hat{\ell}_0 = \frac{\ell_t - \sqrt{1-\bar{\alpha}_t} \cdot \epsilon_\theta(\ell_t, t)}{\sqrt{\bar{\alpha}_t}}
\end{equation}

The density is recovered via:
\begin{equation}
\hat{c}(u,v) = \exp(\text{clip}(\hat{\ell}_0, -20, 20))
\end{equation}

\subsubsection{Copula Constraint Enforcement}

To satisfy equations \eqref{eq:normalization}--\eqref{eq:uniform_v}, a Sinkhorn projection is applied:

\begin{algorithm}[H]
\caption{Sinkhorn Projection for Uniform Marginals}
\begin{algorithmic}[1]
\REQUIRE Density grid $c \in \mathbb{R}^{m \times m}$, iterations $K$
\ENSURE Projected density $\tilde{c}$ with uniform marginals
\STATE $\tilde{c} \gets c$
\FOR{$k = 1$ to $K$}
    \STATE $\tilde{c} \gets \tilde{c} \oslash (\text{rowsum}(\tilde{c}) \otimes \mathbf{1}^\top)$ \COMMENT{Normalize rows}
    \STATE $\tilde{c} \gets \tilde{c} \oslash (\mathbf{1} \otimes \text{colsum}(\tilde{c})^\top)$ \COMMENT{Normalize columns}
\ENDFOR
\RETURN $\tilde{c}$
\end{algorithmic}
\end{algorithm}

where $\oslash$ is elementwise division, $\otimes$ is outer product, and rowsum/colsum compute marginal sums.

\subsubsection{Multi-Objective Loss Function}

The total training loss combines multiple terms:
\begin{equation}
\mathcal{L}_{\text{total}} = \mathcal{L}_{\text{noise}} + w_{\text{ISE}} \cdot \mathcal{L}_{\text{ISE}} + w_{\text{tail}} \cdot \mathcal{L}_{\text{tail}} + w_{\text{KL}} \cdot \mathcal{L}_{\text{marg-KL}}
\end{equation}

\paragraph{Integrated Squared Error (ISE) Loss:}
\begin{equation}
\mathcal{L}_{\text{ISE}} = \mathbb{E}\left[ \|\log \tilde{c}_\theta - \log c_{\text{target}}\|^2 \right]
\end{equation}
where $\tilde{c}_\theta$ is the projected reconstructed density.

\paragraph{Tail Loss:}
Samples points $(u_i, v_i)$ uniformly from corner regions $[0, \tau] \times [0, \tau]$ and $[1-\tau, 1] \times [1-\tau, 1]$:
\begin{equation}
\mathcal{L}_{\text{tail}} = -\mathbb{E}_{(u,v) \sim \text{corners}} \left[ \log \tilde{c}_\theta(u,v) \right]
\end{equation}

\paragraph{Marginal KL Regularization:}
\begin{equation}
\mathcal{L}_{\text{marg-KL}} = \text{KL}\left( \int \tilde{c}_\theta \, dv \,\|\, \text{Uniform}(0,1) \right) + \text{KL}\left( \int \tilde{c}_\theta \, du \,\|\, \text{Uniform}(0,1) \right)
\end{equation}

\subsection{Implementation Details}

\begin{itemize}
    \item \textbf{Model}: U-Net with base channels 48, channel multipliers $[1, 2, 2]$, attention at resolution 16
    \item \textbf{Grid resolution}: $m = 64$ (discretized $[0,1]^2$)
    \item \textbf{Optimizer}: Adam with learning rate $2 \times 10^{-4}$
    \item \textbf{Batch size}: 8
    \item \textbf{Training steps}: 5,000--10,000
    \item \textbf{Projection iterations}: $K = 15$ (Sinkhorn)
\end{itemize}

\section{Empirical Results and Failure Analysis}

\subsection{Catastrophic Model Collapse}

Across all training runs, the model exhibits complete failure:

\begin{table}[h]
\centering
\begin{tabular}{lcccc}
\hline
\textbf{Run} & \textbf{Config} & \textbf{MAE} & \textbf{MSE} & \textbf{Marginal Dev} \\
\hline
46039309 & Baseline & 0.9997 & 2654 & 0.970 \\
46126299 & Improved projection & 1.0000 & 2654 & 0.000 \\
46127314 & Log-space ISE & 1.0000 & 2654 & 0.000 \\
46128400 & Fixed data norm & 0.9998 & 2654 & 0.000 \\
46149822 & Noise only & 0.9998 & 2447 & 0.000 \\
46162978 & Scaled ISE & 0.9998 & 2447 & 0.000 \\
\hline
\end{tabular}
\caption{Training results showing consistent collapse to uniform density. MAE $\approx 1.0$ indicates the model predicts approximately uniform density everywhere, while ground truth copulas have structured densities with mean value $\approx 1.0$ but large variance.}
\end{table}

\subsection{Loss Progression Analysis}

\subsubsection{Noise-Only Training (Run 46149822)}

Training with $\mathcal{L} = \mathcal{L}_{\text{noise}}$ only:
\begin{itemize}
    \item Noise loss decreases: $1.05 \to 0.02$--$0.10$
    \item Model learns to denoise log-densities
    \item But reconstruction after projection yields uniform density
    \item \textbf{Conclusion}: Noise loss alone insufficient for copula structure
\end{itemize}

\subsubsection{Multi-Objective Training with Original Weights}

Loss weights: $w_{\text{ISE}} = 0.3$, $w_{\text{tail}} = 2.0$

Observed loss magnitudes:
\begin{align}
\mathcal{L}_{\text{noise}} &\approx 0.0001 \\
\mathcal{L}_{\text{ISE}} &\approx 47\text{--}62 \\
\mathcal{L}_{\text{tail}} &\approx 0.0002
\end{align}

Effective contributions:
\begin{align}
\text{Noise contribution} &= 1.0 \times 0.0001 = 0.0001 \\
\text{ISE contribution} &= 0.3 \times 50 = 15 \\
\text{Tail contribution} &= 2.0 \times 0.0002 = 0.0004
\end{align}

\textbf{ISE dominates by factor of $150{,}000\times$}, causing model to ignore diffusion objective.

\subsubsection{Scaled Weights (Run 46162978)}

Loss weights: $w_{\text{ISE}} = 0.0003$, $w_{\text{tail}} = 0.002$

Observed:
\begin{align}
\mathcal{L}_{\text{noise}} &\approx 0.02\text{--}0.19 \\
\mathcal{L}_{\text{ISE}} &\approx 77\text{--}130 \\
\mathcal{L}_{\text{tail}} &\approx 13\text{--}14
\end{align}

Contributions:
\begin{align}
\text{Noise} &\approx 0.03 \\
\text{ISE} &= 0.0003 \times 100 \approx 0.03 \\
\text{Tail} &= 0.002 \times 13 \approx 0.026
\end{align}

Even with extreme scaling, ISE and tail losses remain comparable to noise. More importantly, \textbf{losses do not decrease} -- they oscillate without convergence.

\section{Root Cause Analysis}

\subsection{The Real Problem: Objective Mismatch and Projection}

The catastrophic failure is not primarily due to mathematical impossibility of log-space diffusion, but rather a \textbf{practical incompatibility} arising from:

\subsubsection{Incorrect Likelihood Formulation}

The negative log-likelihood term used was:
\begin{equation}
\mathcal{L}_{\text{NLL}} = -\frac{1}{m^2} \sum_{i,j} \log(D_{\text{pred}}^{ij} + \epsilon)
\end{equation}
which treats all grid cells as equally likely samples. This is minimized (for fixed normalization) by the \emph{uniform density itself}. 

\textbf{What's needed}: Either cross-entropy against analytical target masses $T_{ij}$, or sample-based NLL via bilinear interpolation on actual pseudo-observations $(u_i, v_i)$:
\begin{equation}
\mathcal{L}_{\text{NLL}} = -\frac{1}{N} \sum_{i=1}^N \log c_\theta(u_i, v_i)
\end{equation}

Without correct likelihood guidance, the model has no incentive to place mass in the right locations.

\subsubsection{Projection as Attractor to Uniformity}

Sinkhorn projection (IPFP) computes the KL-projection onto doubly stochastic matrices:
\begin{equation}
\tilde{c} = \arg\min_{c' \in \mathcal{C}} \text{KL}(c' \| c_{\text{pred}})
\end{equation}
where $\mathcal{C}$ is the set of densities with uniform marginals.

If $c_{\text{pred}}$ has no stable mass structure (random or weakly structured), the closest doubly stochastic matrix may be nearly uniform. Combined with gradients flowing through IPFP, the path of least resistance becomes:
\begin{enumerate}
    \item Predict nearly constant log-density
    \item IPFP enforces uniform marginals → near-uniform density
    \item All losses (marginal KL, weak ISE) are satisfied with minimal structure
\end{enumerate}

\textbf{Key insight}: IPFP is an excellent final layer when the network learns densities "almost on the copula manifold," but a dangerous flattening operator when the network outputs unstructured predictions.

\subsubsection{Generative Story vs. Training Story Mismatch}

The diffusion model learns:
\begin{equation}
\text{"Given noisy } \ell_t \text{, predict noise } \epsilon"
\end{equation}
in log-space, but we \emph{never use this generatively}. Instead, we always:
\begin{equation}
\text{Reconstruct } \hat{\ell}_0 \to \text{exponentiate} \to \text{IPFP in density space}
\end{equation}

The gradients flow through the entire chain, but the diffusion objective (learn to reverse Gaussian noise in log-space) is disconnected from the constraint enforcement (IPFP in density space). The model learns "predict something IPFP likes" rather than "learn a rich generative model of log-densities."

\subsection{Log-Space vs Density-Space: A Complication, Not an Impossibility}

\begin{theorem}[Log-Space Diffusion Complicates Copula Constraints]
Let $\ell(u,v) = \log c(u,v)$ be the log-density of a copula. The space of log-densities $\mathcal{L} = \{\ell : \exp(\ell) \text{ is a valid copula}\}$ has the following properties:
\begin{enumerate}
    \item Not closed under Gaussian perturbations
    \item Does not form a linear subspace
    \item Has constraints (uniform marginals) that are \emph{nonlinear} in $\ell$
\end{enumerate}
\end{theorem}

\begin{proof}
\textbf{(1) Not closed:} Consider a valid copula density $c(u,v)$ with log-density $\ell_0$. Add Gaussian noise: $\ell_t = \ell_0 + \epsilon$, $\epsilon \sim \mathcal{N}(0, \sigma^2)$. Then $c_t = \exp(\ell_t) = c_0 \cdot \exp(\epsilon)$ has marginals:
\begin{equation}
\int_0^1 c_t(u,v) \, dv = \int_0^1 c_0(u,v) \exp(\epsilon(u,v)) \, dv
\end{equation}
which is \emph{not} generally equal to 1, violating \eqref{eq:uniform_u}.

\textbf{(2) Not a subspace:} The uniform marginal constraint \eqref{eq:uniform_u} in log-space becomes:
\begin{equation}
\int_0^1 \exp(\ell(u,v)) \, dv = 1
\end{equation}
This constraint is \emph{not linear} in $\ell$.

\textbf{(3) Nonlinear constraints:} For $\ell_1, \ell_2 \in \mathcal{L}$ and $\lambda \in (0,1)$:
\begin{equation}
\int_0^1 \exp(\lambda \ell_1 + (1-\lambda) \ell_2) \, dv \neq 1
\end{equation}
in general, so $\lambda \ell_1 + (1-\lambda) \ell_2 \notin \mathcal{L}$.
\end{proof}

\begin{remark}
This is \emph{not unique to copulas}. In standard image diffusion, the forward process immediately leaves the natural image manifold; the score model learns to map back. The issue here is that we combine:
\begin{itemize}
    \item Noising in log-space (leaving copula manifold)
    \item Constraint enforcement in density-space (IPFP)
    \item Weak likelihood guidance (incorrect NLL)
\end{itemize}
This combination, not log-space diffusion per se, causes collapse.
\end{remark}

\subsection{Unbounded Log-Space Magnitudes}

\begin{lemma}[Log-Space Loss Mismatch]
The ISE loss in log-space measures relative errors:
\begin{equation}
(\log c_1 - \log c_2)^2 = \left(\log \frac{c_1}{c_2}\right)^2
\end{equation}
This over-penalizes relative errors in low-density regions and under-penalizes absolute errors in high-density regions. For copulas, what matters is correct marginal integrals and accurate mass placement, not uniform relative accuracy everywhere.
\end{lemma}

\textbf{Example}: Consider $c_{\text{pred}} = 0.01$ and $c_{\text{target}} = 1000$. Then:
\begin{align}
\text{Density-space error:} \quad &|c_{\text{pred}} - c_{\text{target}}| = 999.99 \\
\text{Log-space error:} \quad &(\log 0.01 - \log 1000)^2 = (-4.6 - 6.9)^2 = 132.25
\end{align}

For $c_{\text{pred}} = 900$ and $c_{\text{target}} = 1000$:
\begin{align}
\text{Density-space error:} \quad &|c_{\text{pred}} - c_{\text{target}}| = 100 \\
\text{Log-space error:} \quad &(\log 900 - \log 1000)^2 = (6.8 - 6.9)^2 = 0.01
\end{align}

\begin{remark}
Log-space losses are useful as secondary stabilizers to prevent extreme values, but should not be the primary training objective for copula density estimation.
\end{remark}

\subsection{Summary of Failure Modes}

The uniform collapse results from a perfect storm:
\begin{enumerate}
    \item \textbf{Wrong objective}: NLL formulation that favors uniform density
    \item \textbf{Weak structure guidance}: No strong sample-based likelihood forcing specific mass placement
    \item \textbf{Projection feedback loop}: IPFP + weak gradients → easiest solution is near-constant prediction
    \item \textbf{Log-space complication}: Copula constraints are nonlinear in log-space, making joint optimization difficult
    \item \textbf{Multi-step overhead}: Full DDPM sampling (50+ steps) inappropriate for per-pair fitting task
\end{enumerate}

\textbf{Conclusion}: The log-space DDPM is not suitable for amortized pair-copula estimation. We need a different approach.

\section{Recommended Solution: Noise-Conditioned Density Denoiser}

For the primary use case (fast pair-copula fitting for vines), we propose treating diffusion as a \textbf{training augmentation technique} rather than a generative sampling process.

\subsection{Core Idea}

\begin{itemize}
    \item \textbf{Training}: Learn to denoise histograms of varying quality (augmented with controlled noise)
    \item \textbf{Inference}: Single forward pass on clean histogram → valid copula density
    \item \textbf{Constraint enforcement}: Always apply Sinkhorn projection (IPFP)
    \item \textbf{Loss}: Direct cross-entropy or sample-based NLL against analytical targets
\end{itemize}

This is conceptually similar to denoising autoencoders, not iterative diffusion sampling.

\subsection{Architecture}

\subsubsection{Input Representation}

Given empirical pseudo-observations $\{(u_i, v_i)\}_{i=1}^N$, construct:
\begin{equation}
H[i,j] = \text{count of samples in cell } (i,j), \quad \text{normalized to } \sum H = 1
\end{equation}

During training, add noise:
\begin{equation}
H_t = \sqrt{\bar{\alpha}_t} H + \sqrt{1-\bar{\alpha}_t} \cdot \eta, \quad \eta \sim \mathcal{N}(0, \mathbf{I})
\end{equation}
where $t \sim \text{Uniform}(0,1)$ and $\bar{\alpha}_t$ follows a cosine schedule.

\textbf{Key difference from DDPM}: $t$ is just an input conditioning variable; we do \emph{not} iteratively denoise at inference time.

\subsubsection{Network Architecture}

2D U-Net taking $(H_t, t)$ as input:
\begin{verbatim}
Input: H_t (1, 64, 64), t (scalar), optional probit coordinates
Architecture:
  - Time embedding: t -> Sinusoidal(128)
  - Coordinate embedding (optional): (u,v) -> (Φ^(-1)(u), Φ^(-1)(v))
  - Encoder:
      Conv(1, 64, 3, stride=1, padding=1) + Time -> GN -> SiLU
      Conv(64, 128, 3, stride=2, padding=1) + Time -> GN -> SiLU  # 64->32
      Conv(128, 256, 3, stride=2, padding=1) + Time -> GN -> SiLU # 32->16
      Conv(256, 512, 3, stride=2, padding=1) + Time -> GN -> SiLU # 16->8
  - Bottleneck:
      Self-Attention(512)
      ResBlock(512, 512) + Time -> GN -> SiLU
  - Decoder:
      ConvTranspose(512, 256, 4, stride=2, padding=1) + skip -> GN -> SiLU
      ConvTranspose(256, 128, 4, stride=2, padding=1) + skip -> GN -> SiLU
      ConvTranspose(128, 64, 4, stride=2, padding=1) + skip -> GN -> SiLU
  - Output head:
      Conv(64, 1, 3, stride=1, padding=1)
      Softplus(·) + ε  # Ensure positivity
\end{verbatim}

Output: $D_{\text{raw}} \in \mathbb{R}^{m \times m}_+$

\subsubsection{Sinkhorn Projection Layer}

Always apply IPFP to enforce copula constraints:
\begin{algorithm}[H]
\caption{Sinkhorn Projection}
\begin{algorithmic}[1]
\REQUIRE $D_{\text{raw}} \in \mathbb{R}^{m \times m}_+$, iterations $K=15$
\STATE $D \gets D_{\text{raw}} / \sum D_{\text{raw}}$ \COMMENT{Normalize to sum=1}
\FOR{$k = 1$ to $K$}
    \STATE $D \gets D \oslash (\text{rowsum}(D) \otimes \mathbf{1}^\top)$ \COMMENT{Uniform row marginals}
    \STATE $D \gets D \oslash (\mathbf{1} \otimes \text{colsum}(D)^\top)$ \COMMENT{Uniform column marginals}
\ENDFOR
\RETURN $\hat{D} = D \cdot m$ \COMMENT{Scale so marginals integrate to 1}
\end{algorithmic}
\end{algorithm}

\subsection{Training Procedure}

\subsubsection{Data Generation}

Synthetic copula zoo with analytical densities:
\begin{itemize}
    \item \textbf{Families}: Gaussian, Clayton, Gumbel, Frank, Joe, Student-$t$
    \item \textbf{Parameters}: Wide range (e.g., $\rho \in [-0.95, 0.95]$, $\theta \in [0.5, 10]$)
    \item \textbf{Mixtures}: Weighted combinations of base copulas
    \item \textbf{Grid}: Compute $c(u,v)$ on $m \times m$ grid (e.g., $m=64$)
\end{itemize}

For each training sample:
\begin{enumerate}
    \item Generate target density $T(u,v)$ analytically
    \item Simulate histogram: sample $N \sim \text{Poisson}(\lambda_N)$ points from $c$, bin into $H$
    \item Add noise: $H_t = \sqrt{\bar{\alpha}_t} H + \sqrt{1-\bar{\alpha}_t} \eta$
    \item Include noise level $t$ as model input
\end{enumerate}

\subsubsection{Loss Function}

\textbf{Primary loss - Cross-Entropy on Grid}:
\begin{equation}
\mathcal{L}_{\text{CE}} = -\sum_{i,j} T_{ij} \log(\hat{D}_{ij} + \epsilon)
\end{equation}
where $T_{ij} = \frac{1}{m^2} c(u_i, v_j)$ and $\sum T_{ij} = 1$.

\textbf{Alternative - Sample-based NLL}:
If training with sampled points $(u_k, v_k) \sim c$:
\begin{equation}
\mathcal{L}_{\text{NLL}} = -\frac{1}{N} \sum_{k=1}^N \log \tilde{c}_\theta(u_k, v_k)
\end{equation}
where $\tilde{c}_\theta$ is obtained via bilinear interpolation on $\hat{D}$.

\textbf{Secondary losses}:
\begin{align}
\mathcal{L}_{\text{ISE}} &= \|\hat{D} - T\|^2 \quad \text{(density-space MSE)} \\
\mathcal{L}_{\text{log-ISE}} &= \alpha \|\log \hat{D} - \log T\|^2 \quad \text{(stabilizer, small } \alpha \text{)} \\
\mathcal{L}_{\text{tail}} &= -\frac{1}{N_{\text{tail}}} \sum_{(u,v) \in \text{corners}} \log \hat{D}(u,v)
\end{align}

\textbf{Total loss}:
\begin{equation}
\mathcal{L} = \mathcal{L}_{\text{CE}} + \lambda_{\text{ISE}} \mathcal{L}_{\text{ISE}} + \lambda_{\text{log}} \mathcal{L}_{\text{log-ISE}} + \lambda_{\text{tail}} \mathcal{L}_{\text{tail}}
\end{equation}

Typical weights: $\lambda_{\text{ISE}} = 0.1$, $\lambda_{\text{log}} = 0.05$, $\lambda_{\text{tail}} = 0.5$.

\subsubsection{Training Algorithm}

\begin{algorithm}[H]
\caption{Train Noise-Conditioned Denoiser}
\begin{algorithmic}[1]
\REQUIRE Dataset $\{(c_i, \text{samples}_i)\}$, network $f_\theta$, epochs $E$
\FOR{epoch = 1 to $E$}
    \FOR{each batch $\{(T_j, H_j)\}$}
        \STATE Sample noise levels: $t_j \sim \text{Uniform}(0, 1)$
        \STATE Sample noise: $\eta_j \sim \mathcal{N}(0, \mathbf{I})$
        \STATE Create noisy histograms: $H_{t,j} = \sqrt{\bar{\alpha}_{t_j}} H_j + \sqrt{1-\bar{\alpha}_{t_j}} \eta_j$
        \STATE Forward pass: $D_{\text{raw},j} = f_\theta(H_{t,j}, t_j)$
        \STATE Project: $\hat{D}_j = \text{Sinkhorn}(D_{\text{raw},j})$
        \STATE Compute loss: $\mathcal{L} = \mathcal{L}_{\text{CE}}(\hat{D}_j, T_j) + \text{secondary terms}$
        \STATE Backpropagate and update $\theta$
    \ENDFOR
    \IF{epoch $\mod 10 = 0$}
        \STATE Validate on held-out copulas, visualize samples
    \ENDIF
\ENDFOR
\end{algorithmic}
\end{algorithm}

\subsection{Inference for Vine Construction}

\begin{algorithm}[H]
\caption{Single-Pass Copula Density Estimation}
\begin{algorithmic}[1]
\REQUIRE Pseudo-observations $\{(u_i, v_i)\}_{i=1}^N$ for a pair, trained model $f_\theta$
\STATE Construct empirical histogram: $H = \text{histogram}(\{(u_i, v_i)\}, \text{bins}=m)$
\STATE Normalize: $H \gets H / \sum H$
\STATE Forward pass: $D_{\text{raw}} = f_\theta(H, t=0)$ \COMMENT{Clean input, no noise}
\STATE Project: $\hat{D} = \text{Sinkhorn}(D_{\text{raw}})$
\STATE Compute $h$-functions via numerical integration:
\begin{align*}
h(u|v) &= \int_0^u \hat{c}(s, v) \, ds \\
h(v|u) &= \int_0^v \hat{c}(u, s) \, ds
\end{align*}
\RETURN Copula density $\hat{c}$ and $h$-functions
\end{algorithmic}
\end{algorithm}

\textbf{Runtime}: One U-Net forward pass ($\sim$5-10ms on GPU) + IPFP ($\sim$1-2ms) = $\sim$10ms per pair.

For a vine with 100 edges: $\sim$1 second total.

\subsection{Advantages for Vine Copulas}

\begin{enumerate}
    \item \textbf{Single-pass inference}: No iterative sampling (50-1000 steps)
    \item \textbf{Exact constraints}: IPFP guarantees uniform marginals
    \item \textbf{Robustness}: Learns to handle varying histogram quality via noise conditioning
    \item \textbf{Correct objective}: CE/NLL directly optimizes likelihood of target copula
    \item \textbf{Scalability}: Fast enough for thousands of pair-copulas in large vines
    \item \textbf{Interpretable}: No latent variables, direct density-space operation
\end{enumerate}

\subsection{Relationship to Diffusion Models}

This approach borrows the noise conditioning strategy from diffusion models:
\begin{itemize}
    \item Training data augmentation: $H_t = \sqrt{\bar{\alpha}_t} H + \sqrt{1-\bar{\alpha}_t} \eta$
    \item Time embedding $t$ conditions the network on noise level
    \item Network learns to "denoise" histograms of varying quality
\end{itemize}

But it \emph{does not} use the generative diffusion process:
\begin{itemize}
    \item No reverse SDE / iterative denoising at inference
    \item No sampling from $\mathcal{N}(0, \mathbf{I})$ to generate novel copulas
    \item $t$ is a conditioning variable, not a time index in a sampling trajectory
\end{itemize}

\textbf{Analogy}: Similar to how BERT uses masked language modeling for pre-training but doesn't iteratively unmask at inference.

\section{Alternative Approaches for Unconditional Generation}

For the separate research goal of unconditional copula generation (Use Case 1), we briefly outline two principled approaches:

\subsection{Latent Diffusion Models}

\subsection{Latent Diffusion Models}

\textbf{Use case}: Unconditional generation of diverse copula families.

\subsubsection{Overview}

Separate representation learning from generation:

\begin{enumerate}
    \item \textbf{Phase 1}: Train autoencoder $E_\phi, D_\psi$ with copula-aware losses
    \item \textbf{Phase 2}: Train DDPM in latent space on $\mathbf{z} = E_\phi(c)$
    \item \textbf{Sampling}: Sample $\mathbf{z} \sim p_\theta(\mathbf{z})$, decode $\hat{c} = D_\psi(\mathbf{z})$, project
\end{enumerate}

\subsubsection{Autoencoder Design}

\textbf{Encoder} $E_\phi: \mathbb{R}^{m \times m} \to \mathbb{R}^d$:
\begin{verbatim}
Conv2d(1, 64, 3, stride=2)   -> BN -> ReLU  # 64->32
Conv2d(64, 128, 3, stride=2)  -> BN -> ReLU  # 32->16
Conv2d(128, 256, 3, stride=2) -> BN -> ReLU  # 16->8
Conv2d(256, 512, 3, stride=2) -> BN -> ReLU  # 8->4
Flatten -> Linear(512*16, d)
\end{verbatim}

\textbf{Decoder} $D_\psi: \mathbb{R}^d \to \mathbb{R}^{m \times m}$:
\begin{verbatim}
Linear(d, 512*16) -> Reshape(512, 4, 4)
ConvTranspose2d(512, 256, 4, stride=2) -> BN -> ReLU
ConvTranspose2d(256, 128, 4, stride=2) -> BN -> ReLU
ConvTranspose2d(128, 64, 4, stride=2)  -> BN -> ReLU
ConvTranspose2d(64, 1, 4, stride=2)    -> Softplus
Sinkhorn projection
\end{verbatim}

\textbf{Training loss}:
\begin{equation}
\mathcal{L}_{\text{AE}} = \|\tilde{c} - c\|^2 + \alpha \|\log \tilde{c} - \log c\|^2 + \lambda_{\text{marg}} \mathcal{L}_{\text{marginal}} + \lambda_{\text{smooth}} \|\nabla \tilde{c}\|^2
\end{equation}

\subsubsection{Latent Diffusion}

Once autoencoder converges, freeze and train DDPM on latent codes:
\begin{equation}
\mathcal{L}_{\text{latent}} = \mathbb{E}_{t, \mathbf{z}_0, \epsilon} \left[ \|\epsilon - \epsilon_\theta(\mathbf{z}_t, t)\|^2 \right]
\end{equation}

Network: Simple MLP or 1D transformer operating on $\mathbb{R}^d$.

\subsubsection{Evaluation}

\textbf{Pros}:
\begin{itemize}
    \item Separates copula constraints (AE) from generation (diffusion)
    \item Fast sampling ($\sim$20-50 steps in low-dimensional latent space)
    \item Can interpolate between copulas via latent codes
\end{itemize}

\textbf{Cons}:
\begin{itemize}
    \item Two-stage training complexity
    \item AE reconstruction quality limits final sample quality
    \item Overkill for per-pair fitting (Use Case 2)
\end{itemize}

\textbf{Recommendation}: Research direction for copula family generation, not for vine construction.

\subsection{Score-Based Generative Models}

\textbf{Use case}: Unconditional generation with principled density-space operation.

Score-based models \cite{song2020score} learn the \textbf{score function} $s_\theta(c, t) \approx \nabla_c \log p_t(c)$ where $p_t(c)$ is the density at noise level $t$.

\subsubsection{Forward SDE}

Defines stochastic differential equation adding noise gradually:
\begin{equation}
dc = \mathbf{f}(c, t) \, dt + g(t) \, dW_t
\end{equation}

\textbf{Variance Exploding (VE)}:  $dc = \sqrt{d[\sigma^2(t)]/dt} \, dW_t$

\textbf{Variance Preserving (VP)}:  $dc = -\frac{1}{2} \beta(t) c \, dt + \sqrt{\beta(t)} \, dW_t$

\subsubsection{Reverse SDE with Score}

Sample by reversing the process using learned score:
\begin{equation}
dc = \left[ \mathbf{f}(c, t) - g(t)^2 \nabla_c \log p_t(c) \right] dt + g(t) \, d\bar{W}_t
\end{equation}

\subsubsection{Training}

Denoising score matching in density space:
\begin{equation}
\mathcal{L}_{\text{DSM}} = \mathbb{E}_{t, c_0, \epsilon} \left[ \left\| s_\theta(c_0 + \sigma(t) \epsilon, t) + \frac{\epsilon}{\sigma(t)} \right\|^2 \right]
\end{equation}

Optional marginal penalty to encourage copula constraints:
\begin{equation}
\mathcal{L} = \mathcal{L}_{\text{DSM}} + \lambda_{\text{marg}} \left( \left\| \int c_t \, dv - 1 \right\|^2 + \left\| \int c_t \, du - 1 \right\|^2 \right)
\end{equation}

\subsubsection{Sampling}

Predictor-corrector algorithm with Sinkhorn projection at each step (Algorithm 5):
\begin{enumerate}
    \item Initialize $c_T \sim \mathcal{N}(\mathbf{1}, \sigma^2_{\max} \mathbf{I})$, project
    \item For $t = T \to 0$:
    \begin{itemize}
        \item Predictor: SDE step using score
        \item Project with Sinkhorn
        \item Corrector: Langevin MCMC refinement
        \item Project again
    \end{itemize}
\end{enumerate}

Requires $\sim$1000 steps, each with U-Net forward pass + IPFP.

\subsubsection{Evaluation}

\textbf{Pros}:
\begin{itemize}
    \item Natural density-space operation
    \item Integrated manifold projection
    \item Exact likelihood via probability flow ODE
    \item Strong theoretical foundation
\end{itemize}

\textbf{Cons}:
\begin{itemize}
    \item Very expensive sampling (1000 steps × IPFP)
    \item Many hyperparameters to tune
    \item Overkill for per-pair vine fitting
\end{itemize}

\textbf{Recommendation}: Principled research direction for unconditional generation; too slow for vine construction.

\section{Comparison and Recommendations}

\subsection{Summary Table}

\begin{table}[H]
\centering
\small
\begin{tabular}{lp{2.8cm}p{2.8cm}p{2.8cm}p{2.8cm}}
\hline
\textbf{Criterion} & \textbf{Log-DDPM (Failed)} & \textbf{Denoiser (Recommended)} & \textbf{Latent Diff} & \textbf{Score SDE} \\
\hline
Use case & Unconditional generation & Per-pair vine fitting & Unconditional generation & Unconditional generation \\
\hline
Training space & Log-density & Density & Latent codes & Density \\
\hline
Constraints & Post-hoc (fails) & IPFP (always) & Decoder + IPFP & Manifold projection \\
\hline
Inference & 50 steps & 1 step & 20 steps & 1000 steps \\
\hline
Runtime/pair & $\sim$50ms & $\sim$10ms & N/A & N/A \\
\hline
Training complexity & Medium & Medium & High (2-phase) & Medium-High \\
\hline
Sample quality & Failed (uniform) & Expected good & Unknown & Expected good \\
\hline
Likelihood eval & Hard & Easy (CE) & Very hard & Easy (ODE) \\
\hline
Implementation & Exists (failed) & Minor changes & New arch & Moderate changes \\
\hline
Status & Abandoned & **Recommended** & Future work & Future work \\
\hline
\end{tabular}
\caption{Comparison of approaches for copula density estimation. The noise-conditioned denoiser is recommended for vine copula construction (Use Case 2), while latent diffusion and score-based models are research directions for unconditional generation (Use Case 1).}
\end{table}

\subsection{Primary Recommendation: Noise-Conditioned Denoiser}

For the main goal of fast, robust pair-copula estimation in vine construction:

\textbf{Implement the noise-conditioned density denoiser (Section 5)}:
\begin{itemize}
    \item Single-pass inference: histogram → U-Net → IPFP → density ($\sim$10ms)
    \item Correct likelihood objective: cross-entropy or sample-based NLL
    \item Robustness via noise augmentation during training
    \item Exact copula constraints via Sinkhorn
    \item Scalable to thousands of pairs
\end{itemize}

\textbf{Expected timeline}: 2-3 weeks
\begin{itemize}
    \item Week 1: Implement corrected loss function (CE/NLL), retrain existing U-Net
    \item Week 2: Validate on copula zoo, tune hyperparameters
    \item Week 3: Integration with vine construction pipeline
\end{itemize}

\textbf{Key changes from failed DDPM}:
\begin{enumerate}
    \item \textbf{Loss}: Replace incorrect NLL with proper CE on grid or sample-based NLL
    \item \textbf{Inference}: Remove iterative denoising, use single forward pass at $t=0$
    \item \textbf{Training}: Keep noise conditioning for robustness, but treat as augmentation
    \item \textbf{Space}: Work in density space throughout (no log-density diffusion)
\end{enumerate}

\subsection{Future Research Directions}

\paragraph{Latent Diffusion (Section 6.1)}
For building a flexible generative model of copula families:
\begin{itemize}
    \item Train copula-aware autoencoder with reconstruction + constraint losses
    \item Train DDPM in latent space on encoded copulas
    \item Applications: copula interpolation, extrapolation, exploration of dependence structures
    \item Timeline: 4-6 weeks for full implementation
\end{itemize}

\paragraph{Score-Based Models (Section 6.2)}
For principled unconditional generation with exact likelihoods:
\begin{itemize}
    \item Learn score function in density space
    \item Reverse SDE with integrated projection
    \item Applications: likelihood-based copula selection, density estimation with uncertainty
    \item Timeline: 4-5 weeks for implementation + tuning
\end{itemize}

Both are viable for Use Case 1 but unnecessary for Use Case 2.

\section{Implementation Checklist}

\subsection{Priority 1: Noise-Conditioned Denoiser (Immediate)}

\subsubsection{Fix Data Pipeline}
\begin{itemize}
    \item[$\checkmark$] Verify densities normalized: $\int c = 1$ (already fixed)
    \item[$\checkmark$] Verify uniform marginals on targets (already validated)
    \item[$\square$] Implement histogram construction from pseudo-observations
    \item[$\square$] Add noise conditioning: $H_t = \sqrt{\bar{\alpha}_t} H + \sqrt{1-\bar{\alpha}_t} \eta$
\end{itemize}

\subsubsection{Fix Loss Function}
\begin{itemize}
    \item[$\square$] **CRITICAL**: Replace uniform-grid NLL with cross-entropy:
    \begin{equation*}
    \mathcal{L}_{\text{CE}} = -\sum_{ij} T_{ij} \log(\hat{D}_{ij} + \epsilon)
    \end{equation*}
    \item[$\square$] Or implement sample-based NLL with bilinear interpolation
    \item[$\square$] Keep ISE, log-ISE, tail as secondary (small weights)
    \item[$\square$] Remove marginal KL (IPFP enforces this)
\end{itemize}

\subsubsection{Modify Inference}
\begin{itemize}
    \item[$\square$] Remove iterative denoising loop
    \item[$\square$] Single forward pass: $f_\theta(H, t=0)$
    \item[$\square$] Always apply Sinkhorn projection
    \item[$\square$] Validate: check marginal uniformity, positivity, normalization
\end{itemize}

\subsubsection{Training}
\begin{itemize}
    \item[$\square$] Train on synthetic copula zoo (existing dataset OK)
    \item[$\square$] Monitor CE loss, ISE, marginal deviation
    \item[$\square$] Visualize predictions every 500 steps
    \item[$\square$] Target metrics: MAE $< 0.3$, MSE $< 500$, ISE $< 100$
\end{itemize}

\subsection{Priority 2: Latent Diffusion (If Needed)}

\subsubsection{Autoencoder}
\begin{itemize}
    \item[$\square$] Design encoder (Conv2d stack to latent dim 256-512)
    \item[$\square$] Design decoder (ConvTranspose2d + Softplus + Sinkhorn)
    \item[$\square$] Train with reconstruction + marginal + smoothness losses
    \item[$\square$] Validate: reconstruction MSE $< 0.05$, preserve structure
\end{itemize}

\subsubsection{Latent DDPM}
\begin{itemize}
    \item[$\square$] Encode dataset to latent codes
    \item[$\square$] Train MLP/transformer in latent space
    \item[$\square$] Sample and decode, validate copula properties
\end{itemize}

\subsection{Priority 3: Score-Based (If Needed)}

\subsubsection{Score Network}
\begin{itemize}
    \item[$\square$] Adapt existing U-Net for score prediction
    \item[$\square$] Train with denoising score matching + marginal penalty
    \item[$\square$] Monitor score norm (should be $O(1)$)
\end{itemize}

\subsubsection{Sampling}
\begin{itemize}
    \item[$\square$] Implement Euler-Maruyama SDE solver
    \item[$\square$] Add Langevin MCMC corrector
    \item[$\square$] Integrate Sinkhorn at each step
    \item[$\square$] Tune: steps, corrector iterations, Langevin rate
\end{itemize}

\section{Conclusion}

We have documented the failure of an unconditional log-space DDPM for copula density estimation, identifying root causes beyond mathematical impossibility:
\begin{itemize}
    \item Incorrect likelihood formulation favoring uniform densities
    \item Objective mismatch between log-space diffusion and density-space constraints
    \item Post-hoc projection creating attractor to uniformity
    \item Wrong tool for the job: full generative model when conditional regression suffices
\end{itemize}

\textbf{The practical solution} for vine copula construction is a noise-conditioned density denoiser that:
\begin{itemize}
    \item Treats noise conditioning as training augmentation, not generative sampling
    \item Uses correct likelihood objectives (cross-entropy or sample-based NLL)
    \item Provides single-pass inference with exact copula constraints
    \item Scales to thousands of pair-copulas efficiently
\end{itemize}

This approach borrows the robustness benefits of diffusion training without the computational overhead of iterative sampling.

For unconditional copula generation research, latent diffusion and score-based models remain viable directions, but they address a fundamentally different problem than fast per-pair estimation for vines.

\textbf{Next steps}: Implement corrected loss function, validate on copula zoo, integrate with vine pipeline. Expected improvement: MAE $1.0 \to 0.2$, enabling practical deep learning for vine copulas.

\bibliographystyle{plain}
\begin{thebibliography}{9}

\bibitem{ho2020denoising}
Jonathan Ho, Ajay Jain, and Pieter Abbeel.
\newblock Denoising diffusion probabilistic models.
\newblock In \emph{NeurIPS}, 2020.

\bibitem{song2020score}
Yang Song, Jascha Sohl-Dickstein, Diederik P. Kingma, Abhishek Kumar, Stefano Ermon, and Ben Poole.
\newblock Score-based generative modeling through stochastic differential equations.
\newblock In \emph{ICLR}, 2021.

\bibitem{rombach2022high}
Robin Rombach, Andreas Blattmann, Dominik Lorenz, Patrick Esser, and Björn Ommer.
\newblock High-resolution image synthesis with latent diffusion models.
\newblock In \emph{CVPR}, 2022.

\bibitem{nelsen2006introduction}
Roger B. Nelsen.
\newblock \emph{An Introduction to Copulas}.
\newblock Springer, 2nd edition, 2006.

\bibitem{cuturi2013sinkhorn}
Marco Cuturi.
\newblock Sinkhorn distances: Lightspeed computation of optimal transport.
\newblock In \emph{NeurIPS}, 2013.

\end{thebibliography}

\end{document}
